\subsection{Self-Supervised Monocular Depth Estimation}

Motivated by the prohibitive annotation cost of depth datasets and inspired by stereo-supervised approaches~\cite{garg2016unsupervised, 8100182}, self-supervised monocular depth estimation (SMD) has emerged as a compelling alternative to supervised learning. The seminal work of Zhou et al.~\cite{zhou2017unsupervised} established the foundational framework by jointly training depth and pose networks to minimize a photometric view-synthesis loss, reconstructing each target frame from adjacent source frames under geometric warping. This formulation relies on two assumptions---scene staticity and the absence of occlusions---that are frequently violated in practice, producing artifacts around moving objects and disoccluded regions. Subsequent work introduced a range of mitigations: minimum-reprojection and auto-masking heuristics to discard unreliable pixels~\cite{godard2019digging}, and auxiliary signals such as optical flow~\cite{casser2019depth, chen2019self, hui2022rm, nguyen2024mining, sun2024dynamo}, scene flow~\cite{casser2019depth}, and semantic segmentation~\cite{klingner2020self} to reason about dynamic motion explicitly. Architectural advances---depth-aware convolutions~\cite{arch1, arch2} and transformer backbones~\cite{monovit, arch3, wang2024sqldepth}---further improved accuracy by modeling long-range context. A number of methods additionally aggregate several frames at test time, typically through cost volumes that encode cross-frame correspondence~\cite{watson2021temporal, woo2024prodepth}; while accurate, these models add computational cost and, like all SMD methods, remain bound by two structural limitations that motivate our work: the photometric objective is scale-ambiguous, so it cannot recover metric depth, and networks trained from scratch on a single dataset generalize poorly beyond their training distribution.

\subsection{Depth Foundation Models}

The success of large-scale pretraining in vision and language has inspired depth foundation models that learn generalizable depth priors from massive, heterogeneous data. Pioneering works such as MegaDepth~\cite{li2018megadepth} and MiDaS~\cite{ranftl2020towards} showed that training on diverse image collections with scale- and shift-invariant (SSI) losses yields impressive zero-shot prediction across unseen domains, and DepthAnything and its successor DepthAnything\,V2~\cite{yang2024depth} pushed affine-invariant estimation to state-of-the-art fidelity in capturing edges and fine geometric detail. Crucially, the SSI training objective is what enables this fidelity: by never committing to an absolute scale, the model is free to allocate capacity to relative structure. The same objective, however, leaves predictions without metric scale, limiting their direct use in navigation and manipulation. Metric foundation models address this by predicting absolute depth---ZoeDepth~\cite{bhat2023zoedepth} via relative-to-metric heads, Metric3D~\cite{yin2023metric3d} and its successor~\cite{liu2024metricdepth} by conditioning on camera intrinsics, and UniDepth~\cite{piccinelli2024unidepth} by pursuing universal metric prediction---but recent benchmarking~\cite{li2025benchdepth} reports that enforcing a metric constraint during training consistently degrades geometric detail relative to affine-invariant models, revealing a trade-off between metric consistency and representation richness. Our work aims to obtain the best of both: the geometric fidelity of an affine-invariant model and the metric scale of a metric one.

\subsection{Parameter-Efficient Adaptation and Metric Scale Recovery}

Rather than training a metric model from scratch, we adapt a frozen affine-invariant foundation model, which requires two ingredients: an adaptation mechanism that does not destroy the pretrained priors, and a source of metric scale that does not rely on depth annotations. For the former, parameter-efficient fine-tuning---most prominently Low-Rank Adaptation (LoRA)~\cite{hu2022lora}---updates a small fraction of weights through low-rank residuals, enabling task adaptation while mitigating catastrophic forgetting; we use it to graft self-supervised, metric behavior onto the foundation model while keeping its backbone frozen. For the latter, while metric supervision is classically obtained from LiDAR or ground-truth depth, GPS and other ego-motion sensors provide an inexpensive, widely available metric reference: Chawla et al.~\cite{chawla2021multimodal} align predicted translation magnitude with GPS-derived displacement to resolve the scale ambiguity of self-supervised pose. We build on this idea but apply it in the setting of foundation-model adaptation, and pair it with an explicit \emph{scale--shift decoupling} mechanism so that this weak metric signal can set the global scale of a relative prediction without disturbing its shape---directly targeting the tension between metric scale and geometric fidelity that prior approaches leave unaddressed.

Closest to our goal of annotation-free metric monocular depth for driving are MonoPP~\cite{elazab2025monopp} and its successor Gamma-from-Mono~\cite{elazab2025gamma}, which recover metric scale in closed form from planar-parallax geometry and the camera's known mounting height above the road plane---MonoPP via a multi-frame teacher that distills static-scene scale into a single-image student, Gamma-from-Mono via a single compact network that regresses a dimensionless plane-deviation field (\(\gamma\)) and converts it to metric depth deterministically. Both, however, train lightweight networks from scratch and hinge on a strong, explicit geometric assumption: a visible, locally planar road surface and a known, calibrated camera pose relative to it. Our route to metric scale is comparatively weak---a single GPS-derived displacement ratio with no scene-geometry assumptions---but this is precisely what lets it generalize beyond flat, well-characterized road scenes (we additionally demonstrate zero-shot transfer to Cityscapes, Section~4.5), and what lets us graft it onto a large pretrained backbone instead of learning depth representations from scratch, inheriting the foundation model's edge fidelity and cross-domain transfer that compact from-scratch networks do not exhibit.

Mechanistically, the work closest to our scale--shift decoupling head is RSA~\cite{zeng2024rsa}, which likewise recovers metric depth by regressing a single global affine transformation that is applied on top of a \emph{frozen} relative-depth backbone---the same separation of ``what carries scale'' from ``what carries shape'' that we exploit. We find this convergence telling rather than discouraging: it suggests the decoupling principle is the right one for adapting frozen relative predictors. The two methods differ, however, in exactly the places that matter for an annotation-free pipeline. RSA conditions its affine on a language caption encoded by a frozen CLIP model and, more importantly, \emph{trains} the affine regressor with a supervised L1 loss against ground-truth metric depth---reintroducing the very annotation dependency that motivates working with foundation models in the first place. Our affine is instead conditioned on the depth encoder's own class tokens (no auxiliary modality or caption pipeline) and is learned purely from photometric reprojection and GPS, with no metric labels at any stage---an honest but, we believe, important difference: the same correction mechanism, but one that can be deployed wherever a camera and a coarse displacement signal are available, and the other that still requires a calibration set with ground-truth depth.
